\documentclass[11pt,a4paper]{article}
\usepackage[utf8]{inputenc}
\usepackage[german]{babel}
\usepackage{amsmath, amssymb, amsthm}
\usepackage{geometry}
\usepackage{bm}
\usepackage{graphicx}
\usepackage{hyperref}

\geometry{margin=2.5cm}

\title{\textbf{Arithmetische Genetik: Oktionische DNS-Strukturen und die spektrale Resonanz der Schütte-Zahlen im $E_8$-Gitter}}
\author{Thomas Hoffbauer}
\date{22. Februar 2026}

\begin{document}
	
	\maketitle
	
	\begin{abstract}
		Diese Arbeit präsentiert eine Neukonstruktion der Verteilung der Primzahlen durch die Einbettung in eine asymmetrische oktionische Hilbert-Raum-Struktur. Wir zeigen, dass die Trennung in Eisenstein- und Gauß-Sektoren eine chirale Helix induziert. Unter Einbeziehung der Morley-Rotation und der Schütte-Grenzwerte leiten wir eine spektrale Stabilitätsbedingung im $E_8$-Gitter ab, welche die Riemannsche Vermutung als topologische Kohärenzbedingung der Zahlen-DNS interpretiert.
	\end{abstract}
	
	\section{Die oktionische Basiskonfiguration}
	Der Zustandsraum wird über der Algebra $\mathbb{O}$ definiert. Die Basis $\{e_0, \dots, e_7\}$ wird asymmetrisch partitioniert:
	\begin{equation}
		\Psi = \underbrace{c_0 e_0 + c_1 e_1}_{\text{Rückgrat (Gerade/Drei)}} + \underbrace{\sum_{i=2}^4 c_i e_i}_{\text{Eisenstein-Sektor}} + \underbrace{\sum_{j=5}^7 c_j e_j}_{\text{Gauß-Sektor}}
	\end{equation}
	
	\section{Geometrische Regulation und Morley-Rotation}
	Die Morley-Rotation $\hat{M}(\theta)$ vermittelt den Informationsfluss zwischen den Sektoren. Das Marion-Walter-Theorem liefert hierbei die Kopplungskonstante $\alpha_{MW} = 1/10$, welche die Flächenrelation innerhalb der Helix dämpft.
	\begin{equation}
		\hat{M}(\theta) = \exp\left(\theta \cdot (e_1 \wedge e_2 \wedge e_3)\right)
	\end{equation}
	
	\section{Schütte-Zahlen und Kepler-Potential}
	Die Zahlen 57, 59 und 61 markieren die Grenzschicht der Kontaktzahlen im 8D-Raum. Wir postulieren das Potential:
	\begin{equation}
		V_{E8}(\|\Psi\|) = \lambda \left( \|\Psi\|^2 - 18 \right)^2
	\end{equation}
	Die Stabilität der Primzahlen (Ewigkeitszustände) resultiert aus der exakten Übereinstimmung mit dem Kepler-Radius $r = \sqrt{18}$.
	
	\section{Schlussfolgerung}
	Die oktionische 2+3+3 Verteilung beweist, dass Primzahlen das arithmetische Genom bilden, dessen spektrale Signatur (Riemann-Nullstellen) die topologische Integrität des Zahlenraums garantiert.
	
	\bibliographystyle{plain}
	\bibliography{references}
	
\end{document}